\documentclass{vvsu}

\vvsuyear{2026}

%%%%%%%%%%%%%%%%%%%

\usepackage{graphicx} % для изображений
\usepackage{tabularray} % для таблиц
\usepackage{siunitx} % для обозначений (процент, градус)
\usepackage{csquotes} % для корректной работы с кавычками
\usepackage{listings} % для листингов кода

% Список путей, где будут искаться изображения и файлы
\graphicspath{{images/}}

% Автор документа
\author{А. Е. Филатов}

% Настройка стилей для листингов кода
\input{listing_styles.tex}

%%%%%%%%%%%%%%%%%%%

\begin{document}

% Шапка
\vvsuhead{\linespread{1}\selectfont{}МИНОБРНАУКИ РОССИИ\\
\vspace{10pt}Федеральное государственное бюджетное образовательное учреждение\\
высшего образования\\
\fontsize{13}{13}\selectfont{}<<ВЛАДИВОСТОКСКИЙ ГОСУДАРСТВЕННЫЙ УНИВЕРСИТЕТ>>\\
(ФГБОУ ВО <<ВВГУ>>)\\
\vspace{10pt}\fontsize{12}{12}\selectfont{}ИНСТИТУТ ИНФОРМАЦИОННЫХ ТЕХНОЛОГИЙ И АНАЛИЗА ДАННЫХ\\
КАФЕДРА ИНФОРМАЦИОННЫХ ТЕХНОЛОГИЙ И СИСТЕМ}

% Название отчета
\title{Отчет\\по лабораторной работе №7}
\subtitle{по дисциплине\\<<Информатика и программирование>>}

% Участники работы
\member{Студент\\ гр. БИН-25-3}{А. Е. Филатов}
\member{Ассистент\\ преподавателя}{М.В. Водяницкий}

% Вывод титульника
\maketitle

% Задание
\begin{addition}{Задание}
  Выполнить задания и оформить отчет по стандартам ВВГУ.

  \textit{\textbf{Задание 1.}}
  Имеется список объектов Фонда с указанием уровня угрозы:

 \begin{lstlisting}[basicstyle=\ttfamily\small]
evaluations = [
    {"name": "Agent Cole", "score": 78},
    {"name": "Dr. Weiss", "score": 92},
    {"name": "Technician Moore", "score": 61},
    {"name": "Researcher Lin", "score": 88}
]
\end{lstlisting}
  Используя sorted и лямбда-выражение, отсортируйте объекты по возрастанию уровня угрозы


  \textit{\textbf{Задание 2.}}  
  Дан список сотрудников Фонда с количеством проведенных смен и стоимостью одной смены:
\begin{lstlisting}[basicstyle=\ttfamily\small]
staff_shifts = [
    {"name": "Dr. Shaw", "shift_cost": 120, "shifts": 15},
    {"name": "Agent Torres", "shift_cost": 90, "shifts": 22},
    {"name": "Researcher Hall", "shift_cost": 150, "shifts": 10}
]
\end{lstlisting}
  Используя map и лямбда-выражение, создайте список общей стоимости работы каждого сотрудника
  Затем найдите максимальную стоимость с помощью max

  \textit{\textbf{Задание 3.}}  
  Дан список персонала с уровнем допуска:
\begin{lstlisting}[basicstyle=\ttfamily\small]
personnel = [
    {"name": "Dr. Klein", "clearance": 2},
    {"name": "Agent Brooks", "clearance": 4},
    {"name": "Technician Reed", "clearance": 1}
]
\end{lstlisting}
  Используя map и лямбда-выражение, создайте новый список, где каждому сотруднику добавляется категория допуска:
\begin{vvsu_itemize}
    \item "Restricted" - уровень 1
    \item "Confidential" - уровни 2–3
    \item "Top Secret" - уровень 4 и выше
\end{vvsu_itemize}
Результат должен быть списком словарей

  \textit{\textbf{Задание 4.}}  
  Дан список зон Фонда с указанием времени активности (в часах):
\begin{lstlisting}[basicstyle=\ttfamily\small]
zones = [
    {"zone": "Sector-12", "active_from": 8, "active_to": 18},
    {"zone": "Deep Storage", "active_from": 0, "active_to": 24},
    {"zone": "Research Wing", "active_from": 9, "active_to": 17}
]
\end{lstlisting}
  Используя filter и лямбда-выражение, выберите зоны, которые полностью работают в дневной период (с 8 до 18 включительно)

  \textit{\textbf{Задание 5.}}  
  Фонд анализирует служебные отчеты. Некоторые отчеты содержат внешние ссылки, которые должны быть удалены перед архивированием
  \begin{lstlisting}[basicstyle=\ttfamily\small]
  reports = [
    {"author": "Dr. Moss", "text": "Analysis completed. Reference: http://external-archive.net"},
    {"author": "Agent Lee", "text": "Incident resolved without escalation."},
    {"author": "Dr. Patel", "text": "Supplementary data available at https://secure-research.org"},
    {"author": "Supervisor Kane", "text": "No anomalies detected during inspection."},
    {"author": "Researcher Bloom", "text": "Extended observations uploaded to http://research-notes.lab"},
    {"author": "Agent Novak", "text": "Perimeter secured. No external interference observed."},
    {"author": "Dr. Hargreeve", "text": "Full containment log stored at https://internal-db.scp"},
    {"author": "Technician Moore", "text": "Routine maintenance completed successfully."},
    {"author": "Dr. Alvarez", "text": "Cross-reference materials: http://crosslink.foundation"},
    {"author": "Security Officer Tan", "text": "Shift completed without incidents."},
    {"author": "Analyst Wright", "text": "Statistical model published at https://analysis-hub.org"},
    {"author": "Dr. Kowalski", "text": "Behavioral deviations documented internally."},
    {"author": "Agent Fischer", "text": "Additional footage archived: http://video-storage.sec"},
    {"author": "Senior Researcher Hall", "text": "All test results verified and approved."},
    {"author": "Operations Lead Grant", "text": "Emergency protocol draft shared via https://ops-share.scp"}
]
\end{lstlisting}
Используя filter и лямбда-выражение:
\begin{vvsu_list}
  \item Отберите отчеты, содержащие ссылки (http или https)
  \item Преобразуйте их так, чтобы вместо ссылки отображалось [ДАННЫЕ УДАЛЕНЫ]
\end{vvsu_list}

  \textit{\textbf{Задание 6.}}  
  Преобразуйте их так, чтобы вместо ссылки отображалось [ДАННЫЕ УДАЛЕНЫ]
\begin{lstlisting}[basicstyle=\ttfamily\small]
scp_objects = [
    {"scp": "SCP-096", "class": "Euclid"},
    {"scp": "SCP-173", "class": "Euclid"},
    {"scp": "SCP-055", "class": "Keter"},
    {"scp": "SCP-999", "class": "Safe"},
    {"scp": "SCP-3001", "class": "Keter"}
]
\end{lstlisting}
  Используя filter и лямбда-выражение, сформируйте список SCP-объектов, которые требуют усиленных мер содержания
<<<<<<< HEAD
К объектам с усиленными мерами относятся все SCP, класс которых не равен "Safe"

  Результат должен быть списком словарей исходного формата

  \textit{\textbf{Задание 7.}}  
  Дан список инцидентов с количеством задействованного персонала:
\begin{lstlisting}[basicstyle=\ttfamily\small]
incidents = [
=======
  ⚠️ К объектам с усиленными мерами относятся все SCP, класс которых не равен "Safe"

  Результат должен быть списком словарей исходного формата
  \textit{\textbf{Задание 7.}}  
  Дан список инцидентов с количеством задействованного персонала:
\begin{lstlisting}[basicstyle=\ttfamily\small]
  incidents = [
>>>>>>> 8f56f06d11c95807f076f103b3aa6c756ccdd28b
    {"id": 101, "staff": 4},
    {"id": 102, "staff": 12},
    {"id": 103, "staff": 7},
    {"id": 104, "staff": 20}
]
\end{lstlisting}
Используя sorted и лямбда-выражение:
\begin{vvsu_list}
  \item Отсортируйте инциденты по количеству персонала
  \item Оставьте только три наиболее ресурсоемких инцидента
\end{vvsu_list}

  \textit{\textbf{Задание 8.}}  
  Дан список протоколов безопасности и их уровней критичности:
\begin{lstlisting}[basicstyle=\ttfamily\small]
protocols = [
    ("Lockdown", 5),
    ("Evacuation", 4),
    ("Data Wipe", 3),
    ("Routine Scan", 1)
]
\end{lstlisting}

Используя map и лямбда-выражение, создайте новый список строк вида:
\begin{lstlisting}[basicstyle=\ttfamily\small]
"Protocol Lockdown - Criticality 5"
\end{lstlisting}


  \textit{\textbf{Задание 9.}}  
Имеется список смен охраны с указанием длительности (в часах):
 \begin{lstlisting}[basicstyle=\ttfamily\small]
shifts = [6, 12, 8, 24, 10, 4]
\end{lstlisting}
Используя filter и лямбда-выражение, выберите только те смены, которые:
  \begin{vvsu_list}
    \item длятся не менее 8 часов
    \item не превышают 12 часов
\end{vvsu_list}

  \textit{\textbf{Задание 10.}}  
  Дан список сотрудников с результатами психологической оценки (от 0 до 100):
 \begin{lstlisting}[basicstyle=\ttfamily\small]
evaluations = [
    {"name": "Agent Cole", "score": 78},
    {"name": "Dr. Weiss", "score": 92},
    {"name": "Technician Moore", "score": 61},
    {"name": "Researcher Lin", "score": 88}
]
\end{lstlisting}
Используя max и лямбда-выражение, определите сотрудника с наивысшей оценкой

Результатом должно быть имя сотрудника и его балл
\end{addition}

% Содержание
\toc

% Глава - Выполнение работы
\section{Выполнение работы}
\label{sec:execution}

% Подглава - Задание 1
\subsection{Задание 1}
<<<<<<< HEAD
  Требуется из заданного кортежа отсортировать уовни угрозс поощью лямбды функции.
На рисунке \ref{fig:../code/7.1.py} представлен код программы.
=======
  Сначала создадим список lists, заменяем число 3 на 30, а после выводим результат. На рисунке 1 представлен код программы.
>>>>>>> 8f56f06d11c95807f076f103b3aa6c756ccdd28b

\begin{vvsu_figure}{Листинг программы для задания 1}{fig:../code/7.1.py}
  \begin{minipage}{.75\textwidth}
    \lstinputlisting[language=Python,basicstyle=\fontsize{10}{10}\linespread{1}\selectfont\ttfamily]{../code/7.1.py}
  \end{minipage}
\end{vvsu_figure}

\begin{vvsu_list}
<<<<<<< HEAD
  \item Создает список под названием objects, содержащий кортежи
  \item Сортирует список objects по второму элементу каждого кортежа (уровню угрозы)
  \item Начинает цикл по отсортированному списку, где i - текущий кортеж
  \item Выводит форматированную строку с названием и уровнем угрозы
=======
  \item создадим спикок из 10 различных целых чисел;
  \item ищем в списке число 3 и заменяем его на 30;
  \item выводим результат
>>>>>>> 8f56f06d11c95807f076f103b3aa6c756ccdd28b
\end{vvsu_list}

% Подглава - Задание 2
\subsection{Задание 2}
<<<<<<< HEAD
  Дан список сотрудников Фонда проведенными смен и стоимостью одной смены, требуется отсортировать используя map и лямбда-выражение список общей стоимости работы каждого сотрудника и максимальный.
На рисунке \ref{fig:../code/7.2.py} представлен код программы.
=======
 Создаем список из 5 различных чисел, а дальше нам нужно превратить его в список квадратов этих чисел. На рисунке 2 представлен код программы.
>>>>>>> 8f56f06d11c95807f076f103b3aa6c756ccdd28b

\begin{vvsu_figure}{Листинг программы для задания 2}{fig:../code/7.2.py}
  \begin{minipage}{.75\textwidth}
    \lstinputlisting[language=Python,basicstyle=\fontsize{10}{10}\linespread{1}\selectfont\ttfamily]{../code/7.2.py}
  \end{minipage}
\end{vvsu_figure}

\begin{vvsu_list}
<<<<<<< HEAD
  \item Создает список словарей с информацией о сотрудниках
  \item Вычисляет общую стоимость для каждого сотрудника
  \item Находит максимальное значение из списка общих стоимостей
  \item Выводит список общих стоимостей
  \item Выводит максимальную стоимость
\end{vvsu_list} 

% Подглава - Задание 3
\subsection{Задание 3}
  Дан список персонала с уровнем допуска, с помощью map и лямбда-выражение, требуется создать новый список, где каждому сотруднику добавляется категория допуска. После выводим в консоль результат.
На рисунке \ref{fig:../code/7.3.py} представлен код программы.
=======
  \item создаем список из 5 различных чисел;
  \item используем генератор списков для возведения каждого числа в квадрат и выводим результат.
\end{vvsu_list} 
% Подглава - Задание 3
\subsection{Задание 3}
  Нам требуется создать произвольный список числе, а далее программа должна найти наибольшее из чисел списка и разделить его на длину списка. После выводим в консоль результат. На рисунке 3 представлен код программы.
>>>>>>> 8f56f06d11c95807f076f103b3aa6c756ccdd28b

\begin{vvsu_figure}{Листинг программы для задания 3}{fig:../code/7.3.py}
  \begin{minipage}{.75\textwidth}
    \lstinputlisting[language=Python,basicstyle=\fontsize{10}{10}\linespread{1}\selectfont\ttfamily]{../code/7.3.py}
  \end{minipage}
\end{vvsu_figure}

\begin{vvsu_list}
<<<<<<< HEAD
  \item Создается список personnel, содержащий три словаря. Каждый словарь хранит имя сотрудника name и уровень допуска clearance;
  \item map применяет лямбда-функцию к каждому элементу списка personnel;
  \item Добавляет новый ключ "category" с вычисляемым значением на основе clearance;
  \item Выводим итоговый список result.
=======
  \item создаем произвольный список числе;
  \item находим максимальное число из списка, далее делим на длину списка и выводим результат пользователю.
>>>>>>> 8f56f06d11c95807f076f103b3aa6c756ccdd28b
\end{vvsu_list}

% Подглава - Задание 4
\subsection{Задание 4}
<<<<<<< HEAD
  Дан список зон Фонда с указанием времени активности (в часах), требуется выбрать зоны, которые полностью работают в дневной период (с 8 до 18 включительно).
На рисунке \ref{fig:../code/7.4.py} представлен код решения.
=======
  Создаем функцию которая проверяет состоит ли кортеж только из чисел, если да, то сортируем числа от еньшего к большему, если неь, то кортеж остается неизменным. На рисунке 4 представлен код решения.
>>>>>>> 8f56f06d11c95807f076f103b3aa6c756ccdd28b

\begin{vvsu_figure}{Листинг программы для задания 4}{fig:../code/7.4.py}
  \begin{minipage}{.75\textwidth}
    \lstinputlisting[language=Python,basicstyle=\fontsize{10}{10}\linespread{1}\selectfont\ttfamily]{../code/7.4.py}
  \end{minipage}
\end{vvsu_figure}

\begin{vvsu_list}
<<<<<<< HEAD
  \item спиок зон Фонда с указанием времени активности (в часах);
  \item используется filter с лямбда-функцией для отбора зон, которые активны с 8 до 18 часов включительно;
  \item преобразование результата filter в список;
  \item вывод отобранных зон.
=======
  \item объявляем функцию sort с параметром х;
  \item проверка, все ли элементы в x - числа (int или float);
  \item если возвращает значение True, то сортируем кортеж и возвращаем отсортированный кортеж;
  \item если возвращает значение False, то возвращаем кортеж без изменений;
  \item проверка работы функции значением True;
  \item проверка работы функции с значением False.
>>>>>>> 8f56f06d11c95807f076f103b3aa6c756ccdd28b
\end{vvsu_list}

% Подглава - Задание 5
\subsection{Задание 5}
<<<<<<< HEAD
  Используя filter и лямбда-выражение, треуется анализировать служебные отчеты, отберя отчеты содержащие ссылки (http или https) и преобразовать их так чтобы вместо ссылки отображалось [ДАННЫЕ УДАЛЕНЫ].
На рисунке \ref{fig:../code/7.5.py} представлен код программы.
=======
  Создае функцию которая получает на ввод словарь и выведет пользователю товар с максимальной и минимальной ценой. На рисунке 5 представлен код программы.
>>>>>>> 8f56f06d11c95807f076f103b3aa6c756ccdd28b

\begin{vvsu_figure}{Листинг программы для задания 5}{fig:../code/7.5.py}
  \begin{minipage}{.75\textwidth}
    \lstinputlisting[language=Python,basicstyle=\fontsize{10}{10}\linespread{1}\selectfont\ttfamily]{../code/7.5.py}
  \end{minipage}
\end{vvsu_figure}

\begin{vvsu_list}
<<<<<<< HEAD
  \item Создается список reports, содержащий 15 словарей. Каждый словарь представляет отчёт с полями author (автор) и text (текст отчёта). Некоторые тексты содержат ссылки (начинающиеся с http или https).;
  \item Импортируется модуль re для работы с регулярными выражениями.;
  \item находит ключ с максимальным значением в словаре products и ложим в переменную max1;
  \item filter() (строка 21) оставляет только те отчёты, в тексте которых есть подстрока http, а map к отфильрованным отчетам и преобразуется в список filtered reports;
  \item Цикл for проходит по каждому отчёту в filtered reports и выводит автора и изменённый текст отчёта с удалёнными ссылками.
=======
  \item объявляем функцию help с параметром х;
  \item ищем товара с минимальной ценой с помощью функции min с использванием функции для получния значения по ключу и ложим значение в переменную min1;
  \item находит ключ с максимальным значением в словаре products и ложим в переменную max1;
  \item возвращам переменные min1 и max1;
  \item создаем словарь товаров products;
  \item вызываем функцию help с аргументом products и ложим возвращаемые значения.
>>>>>>> 8f56f06d11c95807f076f103b3aa6c756ccdd28b
\end{vvsu_list}

% Подглава - Задание 6
\subsection{Задание 6}
<<<<<<< HEAD
  Дан список SCP-объектов с указанием их класса содержания, требуется сформировать список объяектов, которые требуют усиленных мер содержания.
На рисунке \ref{fig:../code/7.6.py} представлен код программы.
=======
Имеется список произволных элементов, программа на основе данного списка создает словарь, где каждый элемент списка будт ключем и значением. На рисунке 6 представлен код программы.
>>>>>>> 8f56f06d11c95807f076f103b3aa6c756ccdd28b

\begin{vvsu_figure}{Листинг программы для задания 6}{fig:../code/7.6.py}
  \begin{minipage}{.75\textwidth}
    \lstinputlisting[language=Python,basicstyle=\fontsize{10}{10}\linespread{1}\selectfont\ttfamily]{../code/7.6.py}
  \end{minipage}
\end{vvsu_figure}

\begin{vvsu_list}
<<<<<<< HEAD
  \item Создается список scp objects, содержащий пять словарей. Каждый словарь описывает объект SCP с полями scp (идентификатор SCP) и class (класс объекта: "Euclid", "Keter" или "Safe");
  \item filter() применяет лямбда-функцию к каждому элементу списка scp objects;
  \item Лямбда функция проверяет условие;
  \item filter() возвращает только те объекты SCP, класс которых не является "Safe" (т.е. объекты, требующие усиленных мер сдерживания);
  \item Результат преобразуется в список enhanced containment;
  \item Выводит элементы списка enhanced containment с использованием оператора *.
=======
  \item создаем списко list1;
  \item объявляем функцию help с параметром х;
  \item создаем пустой словар dict1;
  \item проодимся списком по всем элементам списка x;
  \item добавление в словарь пары ключ-значение, где и ключ и значение - сам элемент;
  \item возврат полученного словаря;
  \item вызов функции help с аргументом list1 и вывод результата.
>>>>>>> 8f56f06d11c95807f076f103b3aa6c756ccdd28b
\end{vvsu_list}

% Подглава - Задание 7
\subsection{Задание 7}
<<<<<<< HEAD
    Дан список инцидентов с количеством задействованного персонала, используя sorted и лямбда-выражение требуется отсортировать инциденты по кол-ву персоналу и оставить только 3 самых энегоёмкие.
  На рисунке \ref{fig:../code/7.7.py} представлен код программы.
=======
  Программа получает на ввод переменную кортеж, где к английскому слову дан ключ-перевод и с помощью новой созданной функцией eng rus с параметром x- перевод слова и word слово на русском.
  На рисунке 7 представлен код программы.
>>>>>>> 8f56f06d11c95807f076f103b3aa6c756ccdd28b

\begin{vvsu_figure}{Листинг программы для задания 7}{fig:../code/7.7.py}
  \begin{minipage}{.75\textwidth}
    \lstinputlisting[language=Python,basicstyle=\fontsize{10}{10}\linespread{1}\selectfont\ttfamily]{../code/7.7.py}
  \end{minipage}
\end{vvsu_figure}

\begin{vvsu_list}
<<<<<<< HEAD
  \item Создаётся список incidents, содержащий четыре словаря. Каждый словарь описывает инцидент с полями id (идентификатор инцидента) и staff (количество персонала, задействованного в инциденте);
  \item sorted() сортирует список incidents по указанному критерию;
  \item Параметр key=lambda x: x["staff"] указывает, что сортировка выполняется по значению ключа "staff" (количество персонала).;
  \item Параметр reverse=True сортировку по убыванию (от большего к меньшему);
  \item [:3] — оператор среза, который берёт первые три элемента из отсортированного списка (три инцидента с наибольшим количеством персонала);
  \item Результат присваивается переменной top;
  \item Выводит список top.
\end{vvsu_list}

% Подглава - Задание 8
\subsection{Задание 8}
  Дан список протоколов безопасности и их уровней критичности, Используя map и лямбда-выражение, создайте новый список строк вида: "Protocol Lockdown - Criticality 5".
На рисунке \ref{fig:../code/7.8.py} представлен код программы.
ы
=======
  \item создается функция eng rus с параметрами x и word;
  \item Цикл по всем парам ключ значение в словаре x;
  \item Проверка, совпадает ли значение словаря с искомым словом;
  \item Если найдено совпадение, возвращается соответствующий английский ключ;
  \item Если совпадений нет, возвращается None;
  \item Создание словаря английский-русский;
  \item Поиск английского слова для русского 'яблоко'.
\end{vvsu_list}


% Подглава - Задание 8
\subsection{Задание 8}
Создаем маленькую программу для игры "Камень-Ножницы-Бумага-Ящерица-Спок". Где пользователь вводит свой выбор, а программа случайным образом выбирает один из вариантов. После чего сравниваются варианты и определяется победитель.
На рисунке 8 представлен код программы.

>>>>>>> 8f56f06d11c95807f076f103b3aa6c756ccdd28b
\begin{vvsu_figure}{Листинг программы для задания 8}{fig:../code/7.8.py}
  \begin{minipage}{.75\textwidth}
    \lstinputlisting[language=Python,basicstyle=\fontsize{10}{10}\linespread{1}\selectfont\ttfamily]{../code/7.8.py}
  \end{minipage}
\end{vvsu_figure}

\begin{vvsu_list}
<<<<<<< HEAD
  \item Создаётся список protocols;
  \item map() применяет лямбда-функцию к каждому элементу списка protocols;
  \item Лямбда-функция, принимает кортеж protocol и возвращает форматированную строку с названием протокола и его критичностью;
  \item Выводит новый список protocols, содержащий отформатированные строки.
=======
  \item Импорт модуля для случайного выбора;
  \item Пользователь вводит свой выбор;
  \item Объявление функции игры;
  \item Создание правил: ключ побеждает значения в списке;
  \item Компьютер случайно выбирает вариант;
  \item Вывод выбора компьютера;
  \item Проверка ничьи;
  \item Возврат ничьи;
  \item Проверка, есть ли выбор компьютера в списке побеждаемых вариантов пользователя;
  \item Возврат победы пользователя;
  \item Иначе Компьютер выйграл;
  \item Запуск игры и вывод результата.
>>>>>>> 8f56f06d11c95807f076f103b3aa6c756ccdd28b
\end{vvsu_list}


% Подглава - Задание 9
\subsection{Задание 9}
<<<<<<< HEAD
  Имеется список смен охраны с указанием длительности (в часах), требуется Используя filter и лямбда-выражение, выберите только те смены, которые:длятся не менее 8 часов, не превышают 12 часов.
На рисунке \ref{fig:../code/7.9.py} представлен код программы.
=======
  Требуется программа, где в списке она находит перую букву и сооздает новый кортеж, где значени, это буква, а ключи, это слова начинающиеся на данную букву. На рисунке 9 представлен код программы.
>>>>>>> 8f56f06d11c95807f076f103b3aa6c756ccdd28b

\begin{vvsu_figure}{Листинг программы для задания 9}{fig:../code/7.9.py}
  \begin{minipage}{.75\textwidth}
    \lstinputlisting[language=Python,basicstyle=\fontsize{10}{10}\linespread{1}\selectfont\ttfamily]{../code/7.9.py}
  \end{minipage}
\end{vvsu_figure}

\begin{vvsu_list}
<<<<<<< HEAD
  \item Имеется список смен охраны с указанием длительности (в часах);
  \item filter() применяет лямбда-функцию к каждому элементу списка shifts, лямбда-функция проверяет условие: 8 <= x <= 12 (смена длится от 8 до 12 часов включительно), filter() возвращает только те значения из списка, которые удовлетворяют условию;
  \item Выводит новый список shifts, содержащий отфильтрованные смены.
=======
  \item Создание списка слов;
  \item Объявление функции dict с параметром x (список слов);
  \item Создание множества первых букв всех слов: {'я', 'г', 'б', 'к', 'а'};
  \item Создание словаря, где ключи первые буквы, значения списки слов на эту букву;
  \item Вызов функции и сохранение результата;
  \item Вызов функции и вывод результата.
>>>>>>> 8f56f06d11c95807f076f103b3aa6c756ccdd28b
\end{vvsu_list}

% Подглава - Задание 10
\subsection{Задание 10}
<<<<<<< HEAD
  Дан список сотрудников с результатами психологической оценки (от 0 до 100), требуется используя max и лямбда-выражение, определите сотрудника с наивысшей оценкой результатом должно быть имя сотрудника и его балл.
На рисунке \ref{fig:../code/7.11.py} представлен код программы.
=======
 Требуется программа для подсчета из кортежа средней оценки студентов с оценкой и самим баллом а так же поиска студента с наибольшей средней оценкой и вывести его имя и балл . На рисунке 10 представлен код программы.
>>>>>>> 8f56f06d11c95807f076f103b3aa6c756ccdd28b

\begin{vvsu_figure}{Листинг программы для задания 10}{fig:../code/7.11.py}
  \begin{minipage}{.75\textwidth}
    \lstinputlisting[language=Python,basicstyle=\fontsize{10}{10}\linespread{1}\selectfont\ttfamily]{../code/7.11.py}
  \end{minipage}
\end{vvsu_figure}

\begin{vvsu_list}
<<<<<<< HEAD
  \item Дан список сотрудников с результатами психологической оценки (от 0 до 100);
  \item Функция max() находит максимальный элемент в списке evaluations, параметр key=lambda x: x["score"] указывает, что сравнение элементов должно выполняться по значению ключа "score" (баллу), max() возвращает целый словарь (элемент списка), у которого значение score является наибольшим, результат присваивается переменной top;
  \item Выводит информацию о сотруднике с наивысшим баллом, используя f-строку.
\end{vvsu_list}

\end{document}
=======
  \item Создание списка кортежей с именами студентов и их оценками;
  \item Создание пустого словаря для средних баллов;
  \item Цикл по каждому студенту: извлекает имя и список оценок;
  \item Вычисление среднего балла: сумма оценок делится на количество;
  \item Добавление в словарь: имя студента → средний балл;
  \item Поиск студента с максимальным средним баллом;
  \item Получение максимального среднего балла;
  \item Вывод результата.
\end{vvsu_list}
  \end{document}
>>>>>>> 8f56f06d11c95807f076f103b3aa6c756ccdd28b
